Chemistry tables for the dysbiosis → TMAO/LPS → ROS/HIF-1α → aortic ECM cascade

1. Key chemical species

| Species | Formula / Structure | Role in pathway | Notes |
|---------|---------------------|-----------------|-------|
| Choline | \(\mathrm{C_5H_{14}NO^+}\) | Dietary precursor | Converted by gut microbes |
| TMA (trimethylamine) | \(\mathrm{(CH_3)_3N}\) | Intermediate | Volatile, fishy odor |
| TMAO | \(\mathrm{(CH_3)_3NO}\) / \(\mathrm{C_3H_9NO}\) | Key driver | Hepatic oxidation product |
| LPS | Lipid A + core + O-antigen (variable MW ~10–20 kDa) | Endotoxin | Gram-negative outer membrane |
| ROS (\(\mathrm{O_2^{\bullet-}}\), \(\mathrm{H_2O_2}\)) | Superoxide, hydrogen peroxide | Oxidative stress | From NOX / mitochondria |
| HIF-1α | Protein (~826 aa) | Transcription factor | Stabilized by ROS |
| MMP-2 / MMP-9 | Zn-dependent endopeptidases | Matrix proteases | Cleave collagen & elastin |
| Elastin | Cross-linked tropoelastin polymer | Structural ECM | Provides elasticity |
| Collagen I/III | Triple-helix \(\mathrm{(Gly-X-Y)_n}\) | Structural ECM | Tensile strength |
| NADPH / O₂ | Cofactors | FMO3 reaction | Required for TMAO synthesis |

2. Core chemical / biochemical reactions

| Step | Reaction | Enzyme / Catalyst | Location |
|------|----------|-------------------|----------|
| 1 | Choline → TMA + acetaldehyde | CutC/D (glycyl-radical lyase) | Gut lumen (microbiota) |
| 2 | \(\mathrm{(CH_3)_3N + O_2 + NADPH + H^+ \to (CH_3)_3NO + NADP^+ + H_2O}\) | FMO3 (flavin monooxygenase 3) | Hepatocytes |
| 3 | LPS + TLR4/MD-2 → MyD88 → NF-κB + NOX2 activation | TLR4 signaling | Endothelium / VSMCs |
| 4 | \(\mathrm{2O_2^{\bullet-} + 2H^+ \to H_2O_2 + O_2}\) | Superoxide dismutase (SOD) | Cytosol / mitochondria |
| 5 | ROS inhibit PHDs → HIF-1α stabilization | Prolyl-hydroxylase domain enzymes | Cytosol / nucleus |
| 6 | HIF-1α + HIF-1β → HRE binding → ↑MMP transcription | Transcriptional complex | Nucleus |
| 7 | Collagen / Elastin + H₂O → peptide fragments | MMP-2 / MMP-9 (Zn²⁺ active site) | Aortic media / adventitia |

3. Simplified stoichiometric matrix (selected reactions)

Rows = species; columns = reactions (R1–R4). Positive = production, negative = consumption.

| Species     | R1 (CutC) | R2 (FMO3) | R3 (ROS gen.) | R4 (MMP) |
|-------------|-----------|-----------|---------------|----------|
| Choline     | –1        | 0         | 0             | 0        |
| TMA         | +1        | –1        | 0             | 0        |
| TMAO        | 0         | +1        | 0             | 0        |
| O₂          | 0         | –1        | –1 (approx.)  | 0        |
| NADPH       | 0         | –1        | 0             | 0        |
| H₂O₂ / ROS  | 0         | 0         | +1            | 0        |
| Collagen/Elastin | 0    | 0         | 0             | –1       |
| Peptides    | 0         | 0         | 0             | +1       |

4. Aortic ECM composition matrix (approximate mass fractions in healthy media)

| Component          | Approx. % dry weight | Primary function          | Degradation by |
|--------------------|----------------------|---------------------------|----------------|
| Elastin            | 30–50 %             | Elastic recoil            | MMP-2, MMP-9, MMP-12 |
| Collagen I         | 20–30 %             | Tensile strength          | MMP-1, MMP-8, MMP-13 |
| Collagen III       | 10–20 %             | Compliance / fine structure | MMP-1, MMP-3, MMP-9 |
| Proteoglycans      | 5–10 %              | Hydration / signaling     | Various MMPs / ADAMTS |
| VSMC cytoskeleton  | Variable            | Contractile / synthetic phenotype | Apoptosis (caspase) |

These tables summarize the chemical species, reactions, stoichiometry, and matrix composition that underlie the ODE system previously integrated by Runge–Kutta methods.
